\chapter{Estadística para armonizar}\label{cap:estadistica}

\objetivos{
\begin{itemize}
  \item Tratar un conteo de focos como una variable aleatoria y estimar su incertidumbre.
  \item Explicar qué es la sobredispersión y por qué los fuegos la presentan.
  \item Estimar la relación entre dos series con una razón, una recta y una ley de potencia, y
  elegir entre ellas.
  \item Calcular un intervalo de confianza por remuestreo.
  \item Validar un método de calibración separando datos de ajuste y datos de prueba.
\end{itemize}}

\section{Por qué hace falta estadística}

Armonizar es estimar una relación entre dos series de conteos y usarla para predecir. Cada
parte de esa frase es estadística: los conteos fluctúan aunque el fenómeno no cambie; una
relación estimada con datos finitos tiene incertidumbre; y una predicción hay que comprobarla
con datos que no se usaron para construirla. Este capítulo reúne las herramientas mínimas. No
exige conocimientos previos más allá de la media.

\section{Un conteo es una variable aleatoria}

Supón que en una celda, cada agosto, las condiciones son idénticas: mismo clima, mismos
agricultores, mismas costumbres. Aun así, el número de focos de cada agosto no sería el mismo.
Un año la pasada coincide con tres quemas y otro con cinco; un día clave está nublado. El conteo
observado es una realización de un proceso con azar.

\begin{definicion}{variable aleatoria, valor esperado y varianza}
Una \textbf{variable aleatoria} es una cantidad cuyo valor depende del azar. Su \textbf{valor
esperado} $\mu$ es el promedio que se obtendría repitiendo la observación infinitas veces. Su
\textbf{varianza} $\sigma^2$ mide cuánto se dispersan los valores alrededor de $\mu$; su raíz
cuadrada, $\sigma$, es la \textbf{desviación estándar} y tiene las mismas unidades que la
variable.
\end{definicion}

Lo que interesa del fuego no es el conteo de un mes concreto, sino el valor esperado que hay
detrás: la <<actividad típica>>. El conteo es una estimación ruidosa de él.

\subsection{El modelo de Poisson}

El modelo más simple para un conteo de sucesos que ocurren de forma independiente y a un ritmo
medio constante es la \termino{distribución de Poisson}. La probabilidad de observar exactamente
$k$ sucesos cuando se esperan $\mu$ es
\begin{equation}
  P(k) = \frac{\mu^k e^{-\mu}}{k!} .
\end{equation}
Tiene una propiedad que conviene memorizar: \textbf{su varianza es igual a su media},
\begin{equation}\label{eq:poisson-var}
  \sigma^2 = \mu \quad\Longrightarrow\quad \sigma = \sqrt{\mu}.
\end{equation}

\begin{figure}[ht]
\centering
\begin{tikzpicture}
\begin{axis}[width=13cm,height=6.2cm,ybar,bar width=4.2pt,xmin=-0.7,xmax=30.7,ymin=0,ymax=0.25,
  xlabel={Número de focos observado $k$},ylabel={Probabilidad},
  ytick={0,0.1,0.2},yticklabels={0,{0,1},{0,2}},
  legend style={at={(0.97,0.95)},anchor=north east,font=\small,draw=ceniza!40},axis lines*=left,
  declare function={pois(\k,\m)=exp(\k*ln(\m)-\m-lnGamma(\k+1));
    lnGamma(\z)=(\z-0.5)*ln(\z)-\z+0.9189385+1/(12*\z);}]
\addplot[fill=cielo!70,draw=none,samples at={0,...,30}] {pois(x,3)};
\addplot[fill=fuego!70,draw=none,samples at={0,...,30}] {pois(x,15)};
\legend{$\mu = 3$,$\mu = 15$}
\end{axis}
\end{tikzpicture}
\caption{Distribución de Poisson para dos valores esperados. Con $\mu = 3$ no es raro observar
0 o 7; con $\mu = 15$ lo habitual es caer entre 8 y 23.}
\label{fig:poisson}
\end{figure}

De la ecuación~\eqref{eq:poisson-var} sale la regla práctica más útil de todo el capítulo. Si
se observan $N>0$ sucesos bajo Poisson, la estimación de máxima verosimilitud de $\mu$
es $N$ y una aproximación a su error estándar es $\sqrt{N}$. Su tamaño relativo es
\begin{equation}\label{eq:rel}
  \frac{\sigma}{\mu} \approx \frac{\sqrt{N}}{N} = \frac{1}{\sqrt{N}} .
\end{equation}

\begin{ejemplo}{cuánto fiarse de un conteo}
Tres celdas registran en un mes 25, 400 y 10\,000 focos. ¿Cuál es la incertidumbre relativa de
cada conteo según el modelo de Poisson?
\solucion $1/\sqrt{25} = 0{,}20$; $1/\sqrt{400} = 0{,}05$; $1/\sqrt{10000} = 0{,}01$. Es decir,
un 20\,{}\%, un 5\,{}\% y un 1\,{}\%. Con 25 focos, una diferencia del 30\,{}\% entre dos años puede
ser puro azar. Para reducir la incertidumbre a la mitad hacen falta cuatro veces más datos.
\end{ejemplo}

\begin{parareto}
La ecuación~\eqref{eq:rel} gobierna el compromiso central del diseño: \textbf{cuanto más fina la
resolución espacial y temporal, menos focos por casilla y más ruido}. Un calendario semanal por
distrito puede tener casillas con cinco focos, cuyo ruido es del 45\,{}\%. Antes de fijar la
resolución de la aplicación, mira cuántas detecciones caen en una casilla típica. Si son pocas,
agrupa más (meses en lugar de semanas, provincias en lugar de distritos) o muestra la
incertidumbre.
\end{parareto}

\subsection{Sobredispersión}

El modelo de Poisson supone sucesos independientes. Los focos no lo son: un incendio grande
genera decenas de focos de golpe, y un día seco y ventoso dispara muchas quemas a la vez. Los
focos llegan en racimos, y eso hace que los conteos varíen más de lo que predice Poisson.

\begin{definicion}{sobredispersión}
Un conteo está sobredisperso cuando su varianza es mayor que su media. El \emph{índice de
dispersión} $D = \sigma^2/\mu$ vale 1 para Poisson y es mayor que 1 con sobredispersión.
\end{definicion}

\begin{ejemplo}{medir la sobredispersión}
Los conteos de agosto de una celda en diez años fueron: 12, 31, 8, 45, 19, 102, 15, 27, 9, 38.
Calcula la media, la varianza y el índice de dispersión.
\solucion La media es $\bar{x} = 306/10 = 30{,}6$. La varianza muestral,
$s^2 = \frac{1}{n-1}\sum (x_i - \bar{x})^2$, resulta $s^2 = 7074{,}4/9 = 786{,}0$, y la
desviación estándar $s = 28{,}0$. El índice de dispersión es $D = 786{,}0/30{,}6 \approx 26$.
Si el proceso fuera de Poisson, la desviación estándar sería $\sqrt{30{,}6} = 5{,}5$; la real es
cinco veces mayor.
\end{ejemplo}

La regla $1/\sqrt{N}$ es una referencia bajo Poisson, no un límite universal. En un
proceso sobredisperso puede subestimar mucho la variabilidad. Un enfoque es estimar la
variabilidad directamente a partir de los datos, por ejemplo con remuestreo
(sección~\ref{sec:bootstrap}). El otro es cambiar la unidad de conteo por una que agrupe menos:
contar celdas-día con detección puede reducir la multiplicidad, aunque no elimina
dependencias ni garantiza una distribución Poisson.

\section{Describir una distribución asimétrica}

Los conteos del ejemplo anterior tienen un rasgo típico: casi todos son pequeños y uno es
enorme. La distribución es \termino{asimétrica}, con una cola larga hacia la derecha. En esas
condiciones la media y la desviación estándar describen mal lo <<típico>>, porque un solo valor
extremo las arrastra.

\begin{definicion}{mediana y percentiles}
La \textbf{mediana} es el valor que deja la mitad de los datos por debajo y la mitad por encima.
El \textbf{percentil} $p$ es el valor que deja por debajo el $p$ por ciento de los datos; la
mediana es el percentil 50. La distancia entre los percentiles 25 y 75 es el \textbf{rango
intercuartílico}.
\end{definicion}

En el ejemplo, los datos ordenados son 8, 9, 12, 15, 19, 27, 31, 38, 45, 102. La mediana es el
promedio de los dos centrales, $(19+27)/2 = 23$, bastante menor que la media de 30,6. Si el año
de 102 focos hubiera tenido 1000, la media saltaría a 120 y la mediana seguiría en 23.

\begin{figure}[ht]
\centering
\begin{tikzpicture}
\begin{axis}[width=13cm,height=4.2cm,xmin=0,xmax=110,ymin=0,ymax=2,
  xlabel={Focos en agosto},ytick=\empty,axis y line=none,axis x line=bottom]
\addplot[only marks,mark=*,mark size=3pt,fuego] coordinates {(8,1)(9,1)(12,1)(15,1)(19,1)(27,1)(31,1)(38,1)(45,1)(102,1)};
\draw[cielo,very thick] (axis cs:23,0.35)--(axis cs:23,1.65) node[above,font=\footnotesize]{mediana 23};
\draw[bosque,very thick,dashed] (axis cs:30.6,0.35)--(axis cs:30.6,1.25);
\node[bosque,font=\footnotesize,anchor=west] at (axis cs:31.5,0.45) {media 30{,}6};
\node[font=\footnotesize,anchor=south] at (axis cs:102,1.15) {año extremo};
\end{axis}
\end{tikzpicture}
\caption{Los diez conteos del ejemplo. La media se desplaza hacia el valor extremo; la mediana
se queda con el grueso de los datos.}
\label{fig:mediana}
\end{figure}

Para datos de fuego, la mediana y los percentiles son descriptores más honestos que la media y
la desviación estándar. Volverán en el capítulo~\ref{cap:calendario} para definir qué es
<<normal>>.

\subsection{La escala logarítmica}

Otra forma de domar la asimetría es trabajar con el logaritmo de los conteos. El logaritmo
comprime los valores grandes y separa los pequeños: en escala logarítmica, pasar de 10 a 20 es
el mismo paso que de 100 a 200. Eso encaja con el modo en que pensamos el fuego, en términos de
<<el doble>> o <<la mitad>>, y hace que las distribuciones se vuelvan más simétricas. Como el
logaritmo de cero no existe, cuando hay ceros se usa $\log(1 + x)$.

\section{Relacionar dos series}

Llegamos al problema central. Para una región, durante el periodo de solape, se dispone de
parejas de conteos del mismo mes: $x_i$ con MODIS e $y_i$ con VIIRS. Se quiere una regla que
convierta uno en otro. La tabla~\ref{tab:pares} da un conjunto ilustrativo de doce parejas que
servirá de ejemplo en todo el resto del capítulo.

\begin{table}[ht]
\centering
\small
\begin{tabular}{l*{12}{r}r}
\toprule
Mes $i$ & 1 & 2 & 3 & 4 & 5 & 6 & 7 & 8 & 9 & 10 & 11 & 12 & Suma \\
\midrule
MODIS $x_i$ & 4 & 9 & 15 & 22 & 35 & 48 & 60 & 75 & 90 & 110 & 130 & 160 & 758 \\
VIIRS $y_i$ & 14 & 30 & 41 & 70 & 95 & 150 & 170 & 230 & 250 & 330 & 360 & 470 & 2210 \\
$y_i/x_i$ & 3,50 & 3,33 & 2,73 & 3,18 & 2,71 & 3,12 & 2,83 & 3,07 & 2,78 & 3,00 & 2,77 & 2,94 & \\
\bottomrule
\end{tabular}
\caption{Doce parejas de conteos mensuales de una región en el periodo de solape, ordenadas de
menor a mayor actividad (datos ilustrativos).}
\label{tab:pares}
\end{table}

El primer paso siempre es dibujarlas. Un \termino{diagrama de dispersión} con MODIS en el eje
horizontal y VIIRS en el vertical (figura~\ref{fig:dispersion}) muestra de inmediato si la
relación es aproximadamente una recta, si pasa por el origen y si hay puntos anómalos.

\begin{figure}[ht]
\centering
\begin{tikzpicture}
\begin{axis}[width=12.5cm,height=7.4cm,xmin=0,xmax=175,ymin=0,ymax=520,
  xlabel={Focos MODIS en el mes, $x$},ylabel={Focos VIIRS en el mes, $y$},
  legend style={at={(0.03,0.96)},anchor=north west,font=\small,draw=ceniza!40},
  grid=major,grid style={ceniza!15}]
\addplot[only marks,mark=*,mark size=2.6pt,black] coordinates {(4,14)(9,30)(15,41)(22,70)(35,95)(48,150)(60,170)(75,230)(90,250)(110,330)(130,360)(160,470)};
\addplot[fuego,very thick,domain=0:175] {2.916*x};
\addplot[cielo,thick,dashed,domain=0:175] {2.415+2.877*x};
\addplot[ceniza,thin,domain=0:175] {x};
\legend{meses del solape,razón: $y = 2{,}92\,x$,recta: $y = 2{,}4 + 2{,}88\,x$,$y = x$ (sensores iguales)}
\end{axis}
\end{tikzpicture}
\caption{Diagrama de dispersión de las parejas de la tabla~\ref{tab:pares} con dos ajustes. Si
ambos sensores contaran lo mismo, los puntos caerían sobre la línea gris.}
\label{fig:dispersion}
\end{figure}

\subsection{La razón}

El modelo más simple supone proporcionalidad: VIIRS cuenta siempre $r$ veces lo que MODIS,
\begin{equation}
  y \approx r\,x .
\end{equation}
Hay dos maneras naturales de estimar $r$, y no dan lo mismo.

\begin{definicion}{razón de sumas y media de razones}
\[
  r_{\text{sumas}} = \frac{\sum y_i}{\sum x_i},
  \qquad
  r_{\text{medias}} = \frac{1}{n}\sum \frac{y_i}{x_i} .
\]
\end{definicion}

Con la tabla~\ref{tab:pares}: $r_{\text{sumas}} = 2210/758 = 2{,}92$ y $r_{\text{medias}} =
3{,}00$. La diferencia viene de cómo pondera cada una. La media de razones da el mismo peso a
cada mes, de modo que el mes con 4 focos de MODIS, cuya razón de 3,50 es casi puro ruido, cuenta
tanto como el mes con 160. La razón de sumas pondera cada mes por su número de focos: los meses
con mucha actividad, que son los mejor medidos, mandan.

\begin{atencion}[Razones con denominador pequeño]
La razón $y_i/x_i$ de un mes con pocos focos es inestable, y si $x_i = 0$ ni siquiera existe.
Nunca promedies razones mensuales o por celda sin ponderar. Usa la razón de sumas, que es
equivalente a una media de razones ponderada por $x_i$ y que tolera los ceros.
\end{atencion}

\subsection{La recta de mínimos cuadrados}

Un modelo algo más flexible permite un término independiente:
\begin{equation}
  y \approx a + b\,x .
\end{equation}
La \termino{regresión por mínimos cuadrados} elige $a$ y $b$ de modo que la suma de los
cuadrados de los \termino{residuos} (las diferencias entre cada $y_i$ observado y el valor
$a + b x_i$ que predice la recta) sea mínima. La solución es
\begin{equation}\label{eq:mco}
  b = \frac{\sum (x_i - \bar{x})(y_i - \bar{y})}{\sum (x_i - \bar{x})^2},
  \qquad
  a = \bar{y} - b\,\bar{x} .
\end{equation}

Con los datos del ejemplo, $\bar{x} = 63{,}2$, $\bar{y} = 184{,}2$, y se obtiene $b = 2{,}88$ y
$a = 2{,}4$. La pendiente es casi igual a la razón de sumas y el término independiente es
pequeño frente a los valores de $y$: los datos son compatibles con la proporcionalidad.

La calidad del ajuste se resume en el \termino{coeficiente de determinación},
\begin{equation}
  R^2 = 1 - \frac{\sum (y_i - \hat{y}_i)^2}{\sum (y_i - \bar{y})^2},
\end{equation}
donde $\hat{y}_i$ es el valor predicho. $R^2$ es la fracción de la variabilidad de $y$ que la
recta explica; aquí vale 0,996.

\begin{atencion}[Un $R^2$ alto no demuestra que el modelo sea bueno]
Cuando los datos abarcan un intervalo muy amplio (de 4 a 160), casi cualquier relación creciente
da un $R^2$ cercano a 1, porque los pocos puntos con valores grandes dominan la suma. $R^2$
tampoco dice si las predicciones para los meses de poca actividad son buenas. Para juzgar un
modelo hay que mirar los residuos y, sobre todo, validarlo con datos nuevos
(sección~\ref{sec:validacion}).
\end{atencion}

¿Tiene sentido el término independiente $a$? Significaría que, cuando MODIS no detecta nada,
VIIRS detecta $a$ focos. Físicamente es razonable: hay meses con solo quemas pequeñas, invisibles
para MODIS. Pero también implica que, al convertir una serie de MODIS a la escala de VIIRS, a
cada mes sin focos se le asignarían $a$ focos inventados. Si $a$ es pequeño, lo más limpio es
forzar la recta a pasar por el origen; la pendiente óptima es entonces
\begin{equation}
  b_0 = \frac{\sum x_i y_i}{\sum x_i^2},
\end{equation}
que con estos datos vale 2,90.

\subsection{La ley de potencia}

Queda una posibilidad más: que la razón no sea constante, sino que cambie con el nivel de
actividad. Mira la última fila de la tabla~\ref{tab:pares}: las razones tienden a ser mayores en
los meses de poca actividad (3,5; 3,3) que en los de mucha (2,8; 2,9). Tiene explicación física:
en los meses tranquilos predominan las quemas pequeñas, que MODIS no ve; en el pico de la
temporada hay fuegos grandes, que ambos detectan.

Ese comportamiento se captura con una \termino{ley de potencia},
\begin{equation}
  y = \alpha\,x^{\beta},
\end{equation}
que se vuelve una recta al tomar logaritmos:
\begin{equation}
  \log y = \log\alpha + \beta\,\log x .
\end{equation}
Basta aplicar la ecuación~\eqref{eq:mco} a $\log x$ y $\log y$. Con los datos del ejemplo
resulta $\beta = 0{,}96$ y $\alpha = 3{,}51$. Un exponente $\beta = 1$ sería proporcionalidad
exacta; $\beta < 1$ indica que la razón disminuye al aumentar la actividad, como se sospechaba.

\begin{figure}[ht]
\centering
\begin{tikzpicture}
\begin{axis}[width=12.5cm,height=6.6cm,xmode=log,ymode=log,xmin=3,xmax=200,ymin=8,ymax=700,
  xlabel={Focos MODIS, $x$ (escala logarítmica)},ylabel={Focos VIIRS, $y$ (escala logarítmica)},
  xtick={5,10,20,50,100,200},xticklabels={5,10,20,50,100,200},
  ytick={10,20,50,100,200,500},yticklabels={10,20,50,100,200,500},
  legend style={at={(0.03,0.96)},anchor=north west,font=\small,draw=ceniza!40},
  grid=major,grid style={ceniza!15}]
\addplot[only marks,mark=*,mark size=2.6pt,black] coordinates {(4,14)(9,30)(15,41)(22,70)(35,95)(48,150)(60,170)(75,230)(90,250)(110,330)(130,360)(160,470)};
\addplot[violeta,very thick,domain=3:200] {3.51*x^0.956};
\addplot[fuego,thick,dashed,domain=3:200] {2.916*x};
\legend{meses del solape,potencia: $y = 3{,}51\,x^{0{,}96}$,razón: $y = 2{,}92\,x$}
\end{axis}
\end{tikzpicture}
\caption{Las mismas parejas en escala logarítmica. Aquí todos los meses pesan parecido, y se
aprecia que los de poca actividad quedan por encima de la recta de razón constante.}
\label{fig:loglog}
\end{figure}

La escala logarítmica tiene además una ventaja técnica. La dispersión de los conteos crece con
su valor (los meses de 400 focos fluctúan en decenas; los de 10, en unidades), una situación
llamada \termino{heterocedasticidad}. No impide calcular mínimos cuadrados, pero afecta
a su eficiencia y a las fórmulas usuales de incertidumbre. Ajustar logaritmos cambia
la pérdida y puede estabilizar la dispersión; hay que comprobarlo. Exponentiar una
predicción logarítmica no recupera automáticamente la media original. Ambos sensores
tienen error de observación.

\begin{table}[ht]
\centering
\renewcommand{\arraystretch}{1.25}
\begin{tabularx}{\textwidth}{l>{\raggedright\arraybackslash}X>{\raggedright\arraybackslash}X}
\toprule
Modelo & A favor & En contra \\
\midrule
Razón $y = r x$ & un solo parámetro; robusta; tolera ceros; fácil de explicar & supone que la proporción no cambia con la actividad \\
Recta $y = a + bx$ & recoge un nivel de fondo que solo ve VIIRS & inventa focos donde MODIS vio cero; dominada por los meses grandes \\
Potencia $y = \alpha x^\beta$ & recoge una razón que cambia con la actividad; ajuste equilibrado & no admite ceros; dos parámetros que estimar con pocos datos \\
\bottomrule
\end{tabularx}
\caption{Tres modelos para relacionar los conteos de dos sensores.}
\label{tab:modelos}
\end{table}

\begin{parareto}
Empieza con la \textbf{razón de sumas}: es defendible, casi imposible de implementar mal y
sirve de línea base. Después compara. Con los datos reales de Perú
(sección~\ref{sec:peru-modelos}), la ley de potencia reduce el error mensual de conversión de
\PeruWAPE{} a \PeruPotWAPE{} por ciento, y supera también a las razones por trimestre y por
mes. El modelo se elige con validación interna, nunca mirando los años de prueba.
\end{parareto}

\section{Incertidumbre de una estimación: el remuestreo}\label{sec:bootstrap}

La razón $r = 2{,}92$ salió de doce meses. Con otros doce meses habría salido otro número. ¿Cuánto
podría variar? Para la razón de sumas con datos sobredispersos no hay una fórmula sencilla, pero
hay un procedimiento general y fácil de programar.

\begin{definicion}{remuestreo (bootstrap)}
El remuestreo consiste en simular nuevas muestras extrayendo, \emph{con reemplazo}, $n$ elementos
de la muestra original de tamaño $n$; calcular el estadístico de interés en cada muestra
simulada; y usar la dispersión de los valores obtenidos como medida de su incertidumbre.
\end{definicion}

<<Con reemplazo>> significa que un mismo mes puede salir elegido varias veces y otro ninguna. La
idea es que la muestra es la mejor imagen disponible de la población, así que remuestrear la
muestra imita lo que ocurriría al tomar muestras nuevas (figura~\ref{fig:bootstrap}).

\begin{figure}[ht]
\centering
\begin{tikzpicture}[node distance=5mm]
  \node[caja,fill=cieloClaro,text width=3.2cm] (a) {Muestra original\\ 12 parejas $(x_i,y_i)$};
  \node[caja,right=12mm of a,text width=3.8cm] (b) {Sortear 12 parejas\\ con reemplazo};
  \node[caja,right=12mm of b,text width=3.2cm] (c) {Calcular\\ $r^* = \sum y/\sum x$};
  \node[caja,below=11mm of c,fill=bosqueClaro,text width=3.2cm] (d) {Ordenar los $r^*$:\\ percentiles 2,5 y 97,5};
  \draw[flujo] (a)--(b); \draw[flujo] (b)--(c);
  \draw[flujo] (c.north) -- ++(0,0.55) -| (b.north) node[pos=0.25,above,font=\footnotesize]{repetir miles de veces};
  \draw[flujo] (c)--(d);
\end{tikzpicture}
\caption{Esquema del remuestreo para obtener el intervalo de confianza de una razón.}
\label{fig:bootstrap}
\end{figure}

\begin{definicion}{intervalo de confianza}
Un intervalo de confianza del 95\,{}\% para un parámetro es un intervalo construido con un
procedimiento que, repetido con muchas muestras, contendría el valor verdadero en el 95\,{}\% de
los casos. El intervalo percentil de remuestreo usa los percentiles 2,5 y 97,5 de las réplicas;
su cobertura real es aproximada y depende de la muestra y de sus dependencias.
\end{definicion}

\begin{ejemplo}{intervalo de confianza de la razón}
Se aplican 5000 remuestreos a las doce parejas de la tabla~\ref{tab:pares} y se calcula en cada
uno la razón de sumas. Los valores ordenados tienen percentil 2,5 igual a 2,83 y percentil 97,5
igual a 3,00. Interpreta el resultado.
\solucion La razón estimada es 2,92, con un intervalo de confianza del 95\,{}\% de 2,83 a 3,00:
una incertidumbre de alrededor de $\pm 3\,{}\%$. Si un mes de la era MODIS tuvo 80 focos, su
equivalente en escala VIIRS es $80 \times 2{,}92 = 234$, y la incertidumbre \emph{de la
conversión} lo sitúa entre 226 y 240. A eso hay que sumar la incertidumbre propia del conteo de
80 focos: bajo Poisson, una desviación estándar relativa sería aproximadamente el
11\,{}\%. No se suma directamente a la anchura de un intervalo del 95\,{}\%, pues
son medidas con niveles distintos. Los sesgos no modelados tampoco están incluidos.
\end{ejemplo}

\begin{atencion}[Remuestrea unidades independientes]
El remuestreo supone que los elementos sorteados son independientes entre sí. Los focos
individuales no lo son, ni los días consecutivos de un mismo incendio. Remuestrea unidades
grandes: meses, o mejor aún, años completos. Remuestrear focos sueltos produce intervalos
engañosamente estrechos.
\end{atencion}

\section{Validación}\label{sec:validacion}

Un modelo suele parecer mejor en los datos con los que se construyó. La pregunta que importa
es si funciona con datos que no ha visto.

\begin{definicion}{conjunto de ajuste y conjunto de prueba}
El \textbf{conjunto de ajuste} (o de entrenamiento) es la parte de los datos usada para estimar
los parámetros. El \textbf{conjunto de prueba} es una parte reservada, no usada en el ajuste,
con la que se mide el error de las predicciones.
\end{definicion}

Para la armonización hay una partición natural y muy convincente: \textbf{dividir por tiempo}.
Se estima la relación con unos años del solape y se comprueba con otros
(figura~\ref{fig:validacion}). Es exactamente la situación real: aprender en un periodo y
aplicar en otro.

\begin{figure}[ht]
\centering
\begin{tikzpicture}[x=0.9cm,y=0.8cm]
  \foreach \a [count=\i from 0] in {2012,2013,2014,2015,2016,2017,2018,2019,2020}{
    \ifnum\i<6
      \fill[cielo!55] (\i,0) rectangle ++(0.92,1);
    \else
      \fill[fuego!55] (\i,0) rectangle ++(0.92,1);
    \fi
    \node[font=\footnotesize] at (\i+0.46,0.5) {\a};
  }
  \draw[decorate,decoration={brace,amplitude=5pt}] (0,1.2)--(5.92,1.2) node[midway,above=6pt,font=\small,align=center]{ajuste: se estima $r$};
  \draw[decorate,decoration={brace,amplitude=5pt}] (6,1.2)--(8.92,1.2) node[midway,above=6pt,font=\small,align=center]{prueba: se predice VIIRS\\ a partir de MODIS};
  \node[font=\small,align=left,anchor=west] at (9.6,0.5) {y se compara con\\ el VIIRS real};
\end{tikzpicture}
\caption{Validación temporal dentro del periodo de solape.}
\label{fig:validacion}
\end{figure}

El error en el conjunto de prueba se resume con unas pocas medidas. Si $\hat{y}_i$ es la
predicción e $y_i$ el valor real, para $m$ casos de prueba:
\begin{align}
  \text{sesgo medio} &= \frac{1}{m}\sum (\hat{y}_i - y_i), \\
  \text{error absoluto medio (MAE)} &= \frac{1}{m}\sum |\hat{y}_i - y_i|, \\
  \text{error porcentual absoluto medio (MAPE)} &= \frac{100}{m}\sum \frac{|\hat{y}_i - y_i|}{y_i} .
\end{align}
El sesgo dice si el método tiende a pasarse o a quedarse corto; el MAE, cuánto se equivoca en
promedio en unidades del indicador. MAPE es una media de errores relativos: no está
definido si algún $y_i=0$ y puede quedar dominado por conteos pequeños. En ese caso
se conserva MAE y puede usarse WAPE si el total es positivo (capítulo~\ref{cap:aplicacion}).

\begin{ejemplo}{validar la razón}
Usa los ocho primeros meses de la tabla~\ref{tab:pares} para estimar la razón y los cuatro
últimos para probarla.
\solucion Ajuste: $\sum x = 268$, $\sum y = 800$, $r = 800/268 = 2{,}985$. Predicciones para
los meses 9 a 12, $\hat{y} = 2{,}985\,x$:
\begin{center}
\begin{tabular}{lrrrr}
\toprule
Mes & $x$ & $\hat{y}$ & $y$ real & error \\
\midrule
9 & 90 & 269 & 250 & $+19$ \\
10 & 110 & 328 & 330 & $-2$ \\
11 & 130 & 388 & 360 & $+28$ \\
12 & 160 & 478 & 470 & $+8$ \\
\bottomrule
\end{tabular}
\end{center}
Sesgo medio: $(19 - 2 + 28 + 8)/4 = +13$ focos. MAE: $(19+2+28+8)/4 = 14$ focos. MAPE:
$(7{,}6 + 0{,}6 + 7{,}8 + 1{,}7)/4 = 4{,}4$ por ciento. El método acierta con un error del orden del
4\,{}\%, pero con un sesgo positivo: sobrestima los meses más activos. Era previsible, porque se
ajustó con meses de poca actividad, donde la razón es mayor. Es la misma señal que captaba la
ley de potencia.
\end{ejemplo}

\begin{atencion}[Fuga de información]
Si los datos de prueba influyen de cualquier modo en el ajuste (por ejemplo, eligiendo entre
modelos según cuál predice mejor el conjunto de prueba, o calculando una climatología con todos
los años antes de separar), el error medido es optimista. Separa primero y no toques el conjunto
de prueba hasta el final.
\end{atencion}

\begin{parareto}
Una diapositiva con <<calibramos con 2012--2017 y predijimos 2018--2020 con un error del
X\,{}\%>> vale más ante un jurado que cualquier descripción del método. Es la diferencia entre
afirmar que la armonización funciona y demostrarlo. Reserva tiempo para esto.
\end{parareto}

\resumen

\begin{itemize}
  \item Un conteo de focos es una variable aleatoria. Con el modelo de \textbf{Poisson}, la
  escala relativa aproximada de un conteo positivo es $1/\sqrt{N}$; no es universal.
  \item Los focos llegan en racimos y los conteos están \textbf{sobredispersos}: la
  variabilidad puede ser mayor. Contar celdas-día reduce multiplicidad, no todos los sesgos.
  \item Para distribuciones asimétricas, la \textbf{mediana} y los \textbf{percentiles}
  describen mejor lo típico que la media y la desviación estándar. La escala logarítmica ayuda.
  \item La relación entre dos sensores puede modelarse con una \textbf{razón}, una
  \textbf{recta} o una \textbf{ley de potencia}. La razón de sumas es el punto de partida
  recomendable; nunca se promedian razones sin ponderar.
  \item $R^2$ alto no garantiza un buen modelo.
  \item El \textbf{remuestreo} da intervalos de confianza sin fórmulas; hay que remuestrear
  unidades independientes (meses o años).
  \item La \textbf{validación temporal} (ajustar con unos años, probar con otros) es la prueba
  decisiva; se resume con sesgo, MAE y MAPE.
\end{itemize}

\ejercicios

\begin{enumerate}
  \item Una provincia registra 64 celdas-día con fuego en septiembre. Da su incertidumbre
  absoluta y relativa según Poisson. ¿Cuántas harían falta para una incertidumbre relativa del
  5\,{}\%?
  \item Para los conteos 3, 5, 4, 60, 6, 2, 7, 5 calcula media, mediana, varianza e índice de
  dispersión. ¿Qué descriptor usarías para comunicar el valor <<típico>>?
  \item Con las seis primeras parejas de la tabla~\ref{tab:pares}, calcula la razón de sumas y
  la media de razones. Explica la diferencia.
  \item Con las mismas seis parejas, calcula la recta de mínimos cuadrados con la
  ecuación~\eqref{eq:mco}. Interpreta el signo y el tamaño del término independiente.
  \item Un ajuste en escala logarítmica da $\beta = 0{,}80$ y $\alpha = 6{,}0$. Calcula la
  razón $y/x$ que predice para $x = 10$ y para $x = 200$. Describe con palabras qué significa.
  \item Describe, paso a paso, cómo obtendrías por remuestreo el intervalo de confianza de una
  razón VIIRS/MODIS a partir de ocho años de solape con datos mensuales, remuestreando años
  completos. ¿Por qué años y no meses?
  \item Repite el ejemplo de validación ajustando con los meses impares y probando con los
  pares. Compara el sesgo con el del ejemplo y explica la diferencia.
  \item Un equipo prueba cinco modelos, elige el de menor MAPE en los años de prueba y reporta
  ese MAPE como el error de su método. ¿Qué tiene de malo y cómo se corrige?
\end{enumerate}

\lecturas

Para una introducción accesible a regresión, remuestreo y validación, \emph{An Introduction to
Statistical Learning} (James, Witten, Hastie y Tibshirani) está disponible gratuitamente en
línea; sus capítulos 3 y 5 cubren con más profundidad lo tratado aquí.
